\documentclass[12pt]{article}

\usepackage{sbc-template}
\usepackage[brazil]{babel}
\usepackage{graphicx}
\usepackage{times}
\usepackage{url}

\sloppy

\title{A Busca pelo Invisível: Uma Leitura Pessoal sobre a Identidade e a Coragem em ``O Conto da Ilha Desconhecida''}

\author{Yago Melo da Costa}

\begin{document}
\maketitle

\begin{resumo}
\emph{O Conto da Ilha Desconhecida}, de José Saramago (Companhia das Letras, 1997), é lido aqui a partir da consulta ao texto integral em transcrição digital. O palácio do rei é examinado no funcionamento das portas e da cadeia de secretários, e no modo como um pedido sem analogia interrompe esse funcionamento. A tese de um mundo completo é seguida pelas autoridades do enredo, do rei ao capitão do porto, dos geógrafos do sonho aos marinheiros, até a trajetória da mulher da limpeza, que sai do palácio pela porta das decisões e aprende a navegar. O conto recusa a redução da viagem a metáfora interior: a caravela recebe o nome da ilha e parte à procura dela.
\end{resumo}

\begin{abstract}
José Saramago's \emph{The Tale of the Unknown Island} (Companhia das Letras, 1997) is read here through the complete text consulted in a digital transcription. The king's palace is examined in the working of its doors and its chain of secretaries, and in the way a request with no analogue interrupts that working. The thesis of a completed world is traced through the story's authorities, from the king to the captain of the port, from the dream's geographers to the sailors, down to the cleaning woman's trajectory, as she leaves the palace through the door of decisions and learns to sail. The tale refuses to reduce the voyage to an interior metaphor: the caravel receives the island's name and sets out in search of it.
\end{abstract}

\section{O conto em treze páginas}

\noindent
\textbf{Edição utilizada:} A leitura e todas as citações deste artigo baseiam-se no texto integral de \emph{O Conto da Ilha Desconhecida}, de José Saramago, publicado originalmente pela Companhia das Letras em 1997 \cite{saramago1997}. O exemplar consultado é uma transcrição digital de circulação gratuita, cuja paginação não corresponde à de nenhuma edição impressa; por isso, e por terem sido conferidas uma a uma no texto lido, as citações não informam números de página.

\emph{O Conto da Ilha Desconhecida} chegou como segundo livro de 2026, e o modo como é construído me chamou a atenção primeiro: ele é feito quase todo de diálogos curtos, sem aspas, em que cada personagem fala a partir de uma função, o rei das portas e dos obséquios, o capitão do porto, a mulher da limpeza, os marinheiros que não embarcam. Em vez de perguntar o que a ilha \emph{simboliza}, quis acompanhar o que o conto \emph{faz} com o pedido que abre a história, e a resposta que encontrei organiza este artigo: a obra distribui uma única afirmação, a de que o mundo conhecido estaria completo, por toda a cadeia de autoridades do enredo, e responde a ela com prática. O homem pede um barco, obstrui a porta e parte. Neste artigo, desconhecido é o que o mapa não registra e o que só aparece depois que alguém parte.

\section{A porta, a cadeia de secretários e o pedido sem analogia}

O conto abre com uma frase curta: ``Um homem foi bater à porta do rei e disse-lhe, Dá-me um barco''. O pedido não pede dinheiro, título ou condecoração, ``em lugar de pedir, como era o costume de todos, um título, uma condecoração, ou simplesmente dinheiro'', e é justamente por isso que o palácio não o reconhece. O palácio se organiza em portas com funções distintas: o rei ``passava todo o tempo sentado à porta dos obséquios (entenda-se, os obséquios que lhe faziam a ele)'', as petições entram por outra, e há uma terceira, ``a das decisões, que é raro ser usada, mas quando é, é''. Entre a petição e o rei há uma cadeia: o primeiro-secretário chama o segundo, este o terceiro, e assim ``até chegar à mulher da limpeza, a qual, não tendo ninguém em quem mandar, entreabria a porta das petições e perguntava pela frincha, Que é que tu queres''. O sistema é serial, ``de acordo com a pragmática das portas, ali só se podia atender um suplicante de cada vez'', e a linguagem do palácio é a do despacho: a resposta ``sim ou não conforme estivesse a maré''.

O homem interrompe o funcionamento do palácio com o corpo: diante da recusa do rei em vir pessoalmente, ele anuncia ``que não saio daqui até que ele venha'' e ``deitou-se ao comprido no limiar, tapando-se com a manta por causa do frio. Entrar e sair, só por cima dele''. A obstrução funciona porque o atendimento é serial: enquanto houver alguém à espera, mais ninguém é atendido. O conto descreve a burocracia com precisão, e o rei também, porque após a ``ponderação entre os benefícios e os prejuízos'' ele vai à porta ao cabo de três dias, ``para saber o que queria o intrometido que se havia negado a encaminhar o requerimento pelas competentes vias burocráticas''. A instituição cede quando o custo da recusa supera o da concessão, e a escolha da porta também entra na conta: perguntada se deve abrir a porta toda ou só um pouco, a mulher da limpeza indaga ``Toda, ou só um bocadinho'', e o rei, que ``não gostava muito de se expor aos ares da rua'', escolhe ``De par em par'' por vergonha de parecer medroso diante da vizinhança reunida.

Diante do rei, o homem repete o pedido, ``Dá-me um barco'', e recusa-se a justificá-lo além do fim: ``Para ir à procura da ilha desconhecida''. Quando o rei responde ``E tu quem és, para que eu to dê'', ele devolve a pergunta: ``E tu quem és, para que não mo dês''. A troca mostra como o rei entende a posse: ``Sou o rei deste reino, e os barcos do reino pertencem-me todos'', diz o monarca, e ouve de volta que ``mais lhes pertencerás tu a eles do que eles a ti'', porque ``tu, sem eles, és nada, e que eles, sem ti, poderão sempre navegar''. A negociação altera o pedido: o rei concede ``um barco, mas a tripulação terás de arranjá-la tu, os meus marinheiros são-me precisos para as ilhas conhecidas''. A instituição entrega o objeto e retém o trabalho, e a viagem terá de organizar-se a si própria. O cartão de visita, escrito ``sobre o ombro da mulher da limpeza'', carrega ainda uma cláusula de consciência: ``não quero ter remorsos na consciência se as coisas lhe correrem mal''. E há um último detalhe de realismo: os gritos que inclinam o rei vêm de aspirantes que, ``mais para se verem livres dele do que por simpatia solidária'', resolveram ``intervir a favor do homem''. A concessão mantém os cálculos do palácio: o rei cede o barco e retém a tripulação.

\section{A tese do mundo completo e o mapa como prova}

O conto faz a negativa ``Já não há ilhas desconhecidas'' circular por todas as instâncias do enredo, sem atribuí-la a um personagem só. Ao rei, responde-se que ``estão todas nos mapas''; o homem retruca que ``Nos mapas só estão as ilhas conhecidas'' e que ``Se eu to pudesse dizer, então não seria desconhecida''; pressionado sobre a fonte da informação (``A ninguém''), conclui com a formulação que governa o conto: ``Simplesmente por que é impossível que não exista uma ilha desconhecida''. O mapa só existe porque deixa de fora o que ainda não foi mapeado; o conto vive dessa contradição.

A mesma tese aparece na fala de cada autoridade. O capitão do porto repete a tese e se diz mestre do rei no assunto, ``O que ele sabe de ilhas, aprendeu-o comigo'', mas em seguida a desmonta, designando-se ``homem da terra'': ``não ignoro que todas as ilhas, mesmo as conhecidas, são desconhecidas enquanto não desembarcarmos nelas''. A ilha conhecida só se conhece no desembarque. Quando o homem declara que irá à procura de ``uma onde nunca ninguém tenha desembarcado'', o capitão responde com a condição de todo o conto, ``Se chegares'', ao que ouve que ``às vezes naufraga-se pelo caminho'', com o pedido de que se registre nos anais ``que o ponto a que cheguei foi esse''. A resposta do capitão resume a posição: ``Queres dizer que chegar, sempre se chega, não serias quem és se não o soubesses já''.

No sonho, a mesma tese aparece como ciência: ``os geógrafos do rei foram ver nos mapas e declararam que ilhas por conhecer é coisa que se acabou desde há muito tempo''. Na cidade, os marinheiros a repõem como senso comum e recusam embarcar, cansados ``à procura de um impossível, como se ainda estivéssemos no tempo do mar tenebroso''. Na viagem, os passageiros a transformam em violência: quando avistam terra, gritam ``Esta é uma ilha do mapa'' e ameaçam ``matar-te-emos se não nos levares lá''. O conto repete a mesma questão no palácio, no porto, no sonho e entre os passageiros; em todos os casos, o homem não demonstra que a ilha existe, ele vai. A aposta dele é uma paridade de probabilidades diante da água ``cor de jade'' e do céu ``como um incêndio'': ``Tanta como a de ser tenebroso o mar''.

\section{A mulher da limpeza e a decisão de partir}

A mulher da limpeza entra em cena como extremidade da cadeia burocrática, pergunta pela frincha e despacha conforme a maré, e o rei, ao se retirar, deixa-a ``a olhar para ele com cara de caso''. O narrador só explicará o olhar adiante: era ``esse o preciso momento em que ela resolveu ir atrás do homem''. A decisão aparece ligada à porta e ao ofício: ``deu a volta e saiu com o balde e a vassoura por outra porta, a das decisões'', pensando ``que já bastava de uma vida a limpar e a lavar palácios, que tinha chegado a hora de mudar de ofício, que lavar a limpar barcos é que era a sua vocação verdadeira, no mar, ao menos, a água nunca lhe faltaria''.

Ela escolhe o barco pelo gosto antes de qualquer pergunta: escondida ``por trás de uns bidões'', pensa ``Para o meu gosto, aquele'', e, quando o capitão aponta justamente aquele, reivindica: ``é o meu barco, é o meu barco''. O narrador qualifica a reivindicação como ``insólita'' e ``a todos os títulos abusiva'', e a justifica em uma frase: ``o barco era aquele de que ela tinha gostado, simplesmente''. O capitão, então, converte o capricho em princípio: ``Gostar é provavelmente a melhor maneira de ter, ter deve ser a pior maneira de gostar''. O rei possui os barcos por título; a mulher escolhe um barco por gosto. O conto mostra os dois modos de posse e prefere o segundo.

O conto descreve a aprendizagem dela em atos concretos. Enfrenta as gaivotas ``aos guinchos, furiosas, de goela aberta'', e ``redemoinhando a vassoura como se fosse um espadão dos antigos, fez debandar o bando assassino''; varre a coberta; examina ``o estado das costuras'' das velas e pensa que ``As velas são os músculos do barco'', lamentando que ``se não se lhes dá uso regularmente, abrandam, amolecem''. O barco não pode parecer ``um galinheiro'', e a pólvora desmunida não importa, porque ``ir em busca duma ilha desconhecida'' não precisa ser ``uma empresa de guerra'': o que importa é a falta de víveres. Na conversa com o homem, ela formula o currículo da viagem: ``para a navegação só há dois mestres verdadeiros, um que é o mar, o outro que é o barco''; ele completa com ``o céu'', ``os ventos'', ``as nuvens''. Os dois enumeram os mestres da navegação em voz alta, e ambos aprendem como o protagonista aprende a navegar: ``Aprenderei no mar''. O conto ainda anota que ``nem sequer havia sido ainda contratada'': a tripulação nasce por decisão, não por contrato.

A sabedoria dela se mede em provisão e tempo. Ao homem apressado, que teme não encontrar tripulantes, ela responde com o provérbio do avô, ``quem vai ao mar avia-se em terra, e mais não era ele marinheiro'', lembrando que a partida pressupõe abastecimento, estação, maré, gente ao cais. E quando ele desconfia de que ela zomba dele, ela devolve a origem de tudo: ``Nunca me riria de quem me fez sair pela porta das decisões''. A frase liga as duas portas do conto, porque foi o corpo do homem no limiar que a moveu à outra porta. O homem obstrui a porta; a mulher sai pela outra. A caravela passa a ter dois comandos porque os dois gestos se completam.

\section{O sonho da floresta que navega e o nome do barco}

O recrutamento fracassa: os recusados repetem a tese completa e se recusam a ``tirar-se do sossego dos seus lares e da boa vida dos barcos de carreira para se meterem em aventuras oceânicas, à procura de um impossível''. Diante da desistência possível, o homem confessa o limite do próprio argumento, ``Como poderia falar-lhes eu duma ilha desconhecida, se não a conheço'', e sustenta a paridade: ``Tanta como a de ser tenebroso o mar''. Aí começa o diálogo sobre o autoconhecimento. Ele declara o fim da viagem: ``quero encontrar a ilha desconhecida, quero saber quem sou eu quando nela estiver''. Ela contesta: ``Não o sabes''; ele responde: ``Se não sais de ti, não chegas a saber quem és''. Aí ela invoca o filósofo do rei, que ``dizia que todo o homem é uma ilha'', e propõe a formulação que soa a sabedoria interior: ``Que é necessário sair da ilha para ver a ilha, que não nos vemos se não nos saímos de nós''. Quando ela a resume em ``Se não saímos de nós próprios'', ele recusa o resumo: ``Não é a mesma coisa''. O protagonista recusa a redução da viagem a metáfora: o autoconhecimento aconteceria diante do espelho, e o que ele quer é o que verá ``quando nela estiver'', isto é, do outro lado do mar.

O sonho seguinte é a viagem levada às suas contradições e resolvida sem argumento. Nele, os marinheiros que recusaram já estão a bordo; aparecem mulheres ``tantas quantos os marinheiros'', animais de criação, árvores e searas; e os passageiros, que ``Andávamos à procura de um sítio melhor para viver'', exigem desembarque na primeira terra povoada, ameaçando ``matar-te-emos se não nos levares lá'' quando ele quer passar adiante. ``Por si mesma, a caravela virou a proa''; a turba desembarca levando ``os patos, os coelhos e as galinhas'', ``os bois, os burros e os cavalos'', e até as gaivotas ``transportando no bico os seus gaivotinhos''. Ao piloto do sonho resta o que a multidão deixou: ``as árvores, os trigos e as flores'', ``as trepadeiras que se enrolavam nos mastros''; a coberta vira ``um campo lavrado e semeado''; as raízes ``já estão penetrando no cavername''; e ``bastará que o vento sopre nas copas e vá encaminhando a caravela ao seu destino''. A recusa da comunidade termina em floresta que navega: ``É uma floresta que navega e se balanceia sobre as ondas''. Então o sonhador ``trancou a roda do leme e desceu ao campo com a foice na mão'', ceifa ``as primeiras espigas'' e ``viu uma sombra ao lado da sua sombra'', e acorda ``abraçado à mulher da limpeza, e ela a ele, confundidos os corpos, confundidos os beliches''.

Resta o nome. ``Mal o sol acabou de nascer'', os dois ``foram pintar na proa do barco, de um lado e do outro, em letras brancas, o nome que ainda faltava dar à caravela''. E ``Pela hora do meio-dia, com a maré, A Ilha Desconhecida fez-se enfim ao mar, à procura de si mesma''. A caravela recebe o nome do destino e parte. O nome está pintado no barco que vai buscá-lo; a ilha e o instrumento passam a compartilhar o mesmo nome. O rei quer a ilha mapeada antes de conceder o barco (``Também me interessam as desconhecidas quando deixam de o ser''); o homem parte antes de saber onde ela está. O conto não resolve a contradição: deixa o barco com o nome da ilha.

\section{O balde, a vassoura e a minha escolha}

Li este conto num momento em que eu precisava decidir entre permanecer onde estava e sair por uma porta que quase ninguém usava. O que ficou comigo foi o balde e a vassoura nas mãos de quem decide que já bastava de um ofício. A decisão da mulher da limpeza não é um gesto súbito de coragem; é o resultado de uma vida de trabalho que ela decide não continuar. O provérbio do avô, ``quem vai ao mar avia-se em terra'', diz o mesmo de outro modo: a coragem precisa de preparação. A viagem prepara-se com víveres, estação e maré, e a porta das decisões é ``raro ser usada, mas quando é, é''.

Em \emph{A Metamorfose}, Gregor não consegue sair da cama; em Saramago, o homem se deita no limiar e a obstrução rende um barco. As duas instituições processam pessoas por função, e ambas mostram o que acontece quando a função falha: em Kafka, a falha é a doença e a casa se reorganiza para eliminar o improdutivo; em Saramago, a falha é um pedido sem analogia, e o palácio cede um barco sem tripulação. ``Precisamos nos livrar disto aí'' fecha uma porta para dentro; ``Nunca me riria de quem me fez sair pela porta das decisões'' abre uma para fora.

Em \emph{Ideias para adiar o fim do mundo}, Krenak desmonta a completude que os geógrafos do sonho declaram: ``ilhas por conhecer é coisa que se acabou desde há muito tempo''. O ``assalto que fazem à nossa ideia de natureza'' \cite[p. 16]{krenak2019} transforma o mundo mapeado em recurso. Krenak também pede preparação e narrativa, ``contar mais uma história'' \cite[p. 27]{krenak2019}; a mulher da limpeza prepara o barco e pinta o nome na proa. As duas respostas são práticas: nenhuma depende de convencer o rei ou os geógrafos.

Saio do conto com uma regra que uso nas decisões: não confundir o mapa com o mundo e não confundir vontade com partida. Antes de dizer que não há ilha, verifico quem desenhou o mapa; antes de partir, abasteço o barco.

\begin{thebibliography}{99}

\bibitem{saramago1997}
SARAMAGO, José.
\emph{O Conto da Ilha Desconhecida}.
São Paulo: Companhia das Letras, 1997.

\bibitem{kafka2022}
KAFKA, Franz.
\emph{A Metamorfose} [\emph{Die Verwandlung}, 1915].
Tradução de Giovane Rodrigues.
São Paulo: Instituto Mojo, 2022. (Coleção Literatura Livre).
Edição bilíngue Português-Alemão. E-book.
ISBN 978-65-89008-23-1.
Disponível em: \url{https://mojo.org.br/wp-content/uploads/2024/08/006-A-metamorfose-978-65-89008-23-1.pdf}.

\bibitem{krenak2019}
KRENAK, Ailton.
\emph{Ideias para adiar o fim do mundo}.
1ª ed. São Paulo: Companhia das Letras, 2019.
ISBN 978-85-359-3241-6.

\end{thebibliography}

\end{document}